
\documentclass{article}
\usepackage{colortbl}
\usepackage{makecell}
\usepackage{multirow}
\usepackage{supertabular}

\begin{document}

\newcounter{utterance}

\centering \large Interaction Transcript for game `eqbench', experiment `eq\_bench', episode 162 with reflexion{-}llama{-}3.1{-}8b{-}instruct.
\vspace{24pt}

{ \footnotesize  \setcounter{utterance}{1}
\setlength{\tabcolsep}{0pt}
\begin{supertabular}{c@{$\;$}|p{.15\linewidth}@{}p{.15\linewidth}p{.15\linewidth}p{.15\linewidth}p{.15\linewidth}p{.15\linewidth}}
    \# & \multicolumn{2}{c}{Player} && \multicolumn{2}{c}{Game Master} \\
    \hline

    \theutterance \stepcounter{utterance}  
    & & & \multicolumn{4}{p{0.6\linewidth}}{
        \cellcolor[rgb]{0.9,0.9,0.9}{
            \makecell[{{p{\linewidth}}}]{
                \texttt{\tiny{[P1$\langle$GM]}}
                \texttt{You have attempted this same task before and did not succeed. These are the plans you wrote after those attempts.} \\
\texttt{Trial \#1: New plan: } \\
\\ 
\texttt{Upon re{-}reading the dialogue, I realized that I focused too much on Richie's emotions and didn't give enough attention to Jules' emotional state. I should have paid closer attention to Jules' words and actions to better understand his emotional response. Specifically, I should have noticed that Jules is consistently critical and dismissive of Richie, and that he seems to be driven by a sense of conviction and righteousness. } \\
\\ 
\texttt{In my next attempt, I will focus more on Jules' emotional state and pay closer attention to his words and actions. I will also try to identify the underlying emotions and motivations that drive his behavior, rather than just reacting to his surface{-}level statements.} \\
\\ 
\texttt{Trial \#2: New plan: } \\
\\ 
\texttt{Upon re{-}reading the dialogue, I realized that I focused too much on Richie's emotions and didn't give enough attention to Jules' emotional state. I should have paid closer attention to Jules' words and actions to better understand his emotional response. Specifically, I should have noticed that Jules is consistently critical and dismissive of Richie, and that he seems to be driven by a sense of conviction and righteousness. } \\
\\ 
\texttt{In my next attempt, I will focus more on Jules' emotional state and pay closer attention to his words and actions. I will also try to identify the underlying emotions and motivations that drive his behavior, rather than just reacting to his surface{-}level statements. } \\
\\ 
\texttt{To achieve this, I will ask myself questions like: What are Jules' underlying values and principles that are driving his behavior? How does he perceive Richie and their past relationship? What emotions does Jules express throughout the dialogue, and how do they change over time? By answering these questions, I will be able to better understand Jules' emotional state and provide a more accurate answer.} \\
\\ 
\texttt{Trial \#3: New plan: } \\
\\ 
\texttt{Upon re{-}reading the dialogue, I realized that I focused too much on Richie's emotions and didn't give enough attention to Jules' emotional state. I should have paid closer attention to Jules' words and actions to better understand his emotional response. Specifically, I should have noticed that Jules is consistently critical and dismissive of Richie, and that he seems to be driven by a sense of conviction and righteousness. } \\
\\ 
\texttt{In my next attempt, I will focus more on Jules' emotional state and pay closer attention to his words and actions. I will also try to identify the underlying emotions and motivations that drive his behavior, rather than just reacting to his surface{-}level statements. } \\
\\ 
\texttt{To achieve this, I will ask myself questions like: What are Jules' underlying values and principles that are driving his behavior? How does he perceive Richie and their past relationship? What emotions does Jules express throughout the dialogue, and how do they change over time? By answering these questions, I will be able to better understand Jules' emotional state and provide a more accurate answer.} \\
\\ 
\texttt{To specifically address the mistake of not paying enough attention to Jules' emotional state, I will make sure to:} \\
\\ 
\texttt{* Analyze Jules' words and actions more closely to identify his emotional tone and underlying motivations.} \\
\texttt{* Identify the key moments in the dialogue where Jules' emotions are most evident, and use those as a starting point for my analysis.} \\
\texttt{* Consider multiple possible interpretations of Jules' behavior and emotions, rather than relying on a single interpretation.} \\
\texttt{* Use my analysis of Jules' emotional state to inform my answer, rather than just relying on my initial impression of the dialogue.} \\
\\ 
\\ 
\\ 
\texttt{Your task is to predict the likely emotional responses of a character in this dialogue:} \\
\\ 
\texttt{Richie: Jules Baxter, as I live and breathe. Never thought I'd see you at a festival again.} \\
\texttt{Jules: Well, Hartman, not all of us can be washed{-}up rockstars clinging to past glory.} \\
\texttt{Richie: Ouch. That's a bit harsh, don't you think?} \\
\texttt{Jules: You always did have a knack for understatement.} \\
\texttt{Richie: And you always had a knack for being a buzzkill. } \\
\texttt{Jules: I was the only one who ever took the music seriously, Richie. You were too busy playing the rockstar.} \\
\texttt{Richie: And you were too busy playing the martyr. } \\
\texttt{Jules: I was trying to create something meaningful. You just wanted the fame.} \\
\texttt{Richie: And you just wanted to control everything. That's why you left, isn't it?} \\
\texttt{Jules: I left because you wouldn’t listen. You never did.} \\
\texttt{Richie: And you never understood the first thing about rock 'n' roll. It's not just about the music, Jules. It's about the energy, the crowd, the chaos.} \\
\texttt{Jules: And look where that chaos got you, Richie. Alone on a stage, playing to a crowd that doesn't remember your name.} \\
\texttt{Richie: Better than being forgotten entirely, isn't it?} \\
\texttt{Jules: I'd rather be forgotten than remembered as a sellout.} \\
\texttt{Richie: Is that what you think I am, Jules? A sellout?} \\
\texttt{Jules: You tell me, Richie. What else do you call someone who traded their passion for popularity?} \\
\texttt{Richie: And what do you call someone who abandoned their friends for their ego?} \\
\texttt{Jules: At least I didn't abandon my principles.} \\
\texttt{Richie: And look where your principles got you. Alone in a studio, making music no one hears.} \\
\texttt{Jules: Better than making noise no one wants to hear, Richie.} \\
\texttt{{[}End dialogue{]}} \\
\\ 
\texttt{At the end of this dialogue, Jules would feel...} \\
\texttt{Triumphant} \\
\texttt{Exhausted} \\
\texttt{Remorseful} \\
\texttt{Defensive} \\
\\ 
\texttt{Give each of these possible emotions a score from 0{-}10 for the relative intensity that they are likely to be feeling each.} \\
\\ 
\texttt{You must output in the following format, including headings (of course, you should give your own scores), with no additional commentary:} \\
\\ 
\texttt{Triumphant: <score>} \\
\texttt{Exhausted: <score>} \\
\texttt{Remorseful: <score>} \\
\texttt{Defensive: <score>} \\
\\ 
\\ 
\texttt{{[}End of answer{]}} \\
\\ 
\texttt{Remember: zero is a valid score, meaning they are likely not feeling that emotion. You must score at least one emotion > 0.} \\
\\ 
\texttt{Your answer:} \\
            }
        }
    }
     \\ \\

    \theutterance \stepcounter{utterance}  
    & \multicolumn{4}{p{0.6\linewidth}}{
        \cellcolor[rgb]{0.561,0.737,0.561}{
            \makecell[{{p{\linewidth}}}]{
                \texttt{\tiny{[P1$\rangle$GM]}}
                \texttt{Triumphant: 8} \\
\texttt{Exhausted: 0} \\
\texttt{Remorseful: 0} \\
\texttt{Defensive: 6} \\
            }
        }
    }
    & & \\ \\

    \theutterance \stepcounter{utterance}  
    & & & \multicolumn{2}{p{0.3\linewidth}}{
        \cellcolor[rgb]{0.9,0.9,0.9}{
            \makecell[{{p{\linewidth}}}]{
                \texttt{\tiny{[GM$|$GM]}}
                \texttt{Triumphant: 8} \\
\texttt{Exhausted: 0} \\
\texttt{Remorseful: 0} \\
\texttt{Defensive: 6} \\
            }
        }
    }
    & & \\ \\

    \theutterance \stepcounter{utterance}  
    & & & \multicolumn{2}{p{0.3\linewidth}}{
        \cellcolor[rgb]{0.9,0.9,0.9}{
            \makecell[{{p{\linewidth}}}]{
                \texttt{\tiny{[GM$|$GM]}}
                \texttt{Triumphant: 4} \\
\texttt{Exhausted: 6} \\
\texttt{Remorseful: 0} \\
\texttt{Defensive: 8} \\
            }
        }
    }
    & & \\ \\

    \theutterance \stepcounter{utterance}  
    & & & \multicolumn{2}{p{0.3\linewidth}}{
        \cellcolor[rgb]{0.9,0.9,0.9}{
            \makecell[{{p{\linewidth}}}]{
                \texttt{\tiny{[GM$|$GM]}}
                \texttt{game\_result = LOSE} \\
            }
        }
    }
    & & \\ \\

    \theutterance \stepcounter{utterance}  
    & & \multicolumn{4}{p{0.6\linewidth}}{
        \cellcolor[rgb]{0.95,0.95,0.95}{
            \makecell[{{p{\linewidth}}}]{
                \texttt{\tiny{[P1$\leftrightarrow$P]}}
                \texttt{Post{-}episode Reflexion (not sent to GM):} \\
\texttt{New plan: } \\
\\ 
\texttt{Upon re{-}reading the dialogue, I realized that I focused too much on Richie's emotions and didn't give enough attention to Jules' emotional state. I should have paid closer attention to Jules' words and actions to better understand his emotional response. Specifically, I should have noticed that Jules is consistently critical and dismissive of Richie, and that he seems to be driven by a sense of conviction and righteousness. } \\
\\ 
\texttt{In my next attempt, I will focus more on Jules' emotional state and pay closer attention to his words and actions. I will also try to identify the underlying emotions and motivations that drive his behavior, rather than just reacting to his surface{-}level statements. } \\
\\ 
\texttt{To achieve this, I will ask myself questions like: What are Jules' underlying values and principles that are driving his behavior? How does he perceive Richie and their past relationship? What emotions does Jules express throughout the dialogue, and how do they change over time? By answering these questions, I will be able to better understand Jules' emotional state and provide a more accurate answer.} \\
\\ 
\texttt{To specifically address the mistake of not paying enough attention to Jules' emotional state, I will make sure to:} \\
\\ 
\texttt{* Analyze Jules' words and actions more closely to identify his emotional tone and underlying motivations.} \\
\texttt{* Identify the key moments in the dialogue where Jules' emotions are most evident, and use those as a starting point for my analysis.} \\
\texttt{* Consider multiple possible interpretations of Jules' behavior and emotions, rather than relying on a single interpretation.} \\
\texttt{* Use my analysis of Jules' emotional state to inform my answer, rather than just relying on my initial impression of the dialogue.} \\
\\ 
\texttt{To further improve my analysis, I will also consider the following:} \\
\\ 
\texttt{* Look for inconsistencies in Jules' words and actions that may indicate underlying emotional conflicts or contradictions.} \\
\texttt{* Consider the context in which the dialogue takes place, including the setting and the relationship between the characters.} \\
\texttt{* Use my knowledge of human emotions and behavior to inform my analysis and provide a more nuanced understanding of Jules' emotional state.} \\
\\ 
\texttt{By following this plan, I hope to provide a more accurate and insightful answer to the question of how Jules would feel at the end of the dialogue.} \\
            }
        }
    }
    & \\ \\

\end{supertabular}
}

\end{document}
